\clearpage\pagestyle{otMain}\setlength{\otBlocEspacement}{2.5mm}\fontsize{11}{15}\selectfont\corrigesChapitre{7}{Utiliser et comparer les fractions}
\begin{blocLivre}{cor-7-1}\begin{corrige}{1}{Lire sur la demi-droite}\demarche{Compter les pas depuis zéro.}
Chaque pas vaut $\frac14$. A est 3 pas après 0 : $\res{A\left(\frac34\right)}$. B est 3 pas après 1 : $\res{B\left(\frac74\right)}$.
\label{fig-cor-7-1}\figureGeo[corrige]{fd-1}
\end{corrige}\end{blocLivre}
\begin{blocLivre}{cor-7-2}\begin{corrige}{2}{Placer des tiers}\demarche{Avancer d’autant de pas que le numérateur.}
C se place 2 pas après 0 et D cinq pas après 0, c’est-à-dire 2 pas après l’entier 1. Les positions figurent en vert sur le corrigé.
\label{fig-cor-7-2}\figureGeo[corrige]{fd-2}
\end{corrige}\end{blocLivre}
\begin{blocLivre}{cor-7-3}\begin{corrige}{3}{Comparer à l’unité}\demarche{Comparer le numérateur au dénominateur.}
$\res{\frac56<1}$ : il manque un sixième. $\res{\frac76>1}$ : un sixième de trop. $\res{\frac66=1}$ exactement.
\end{corrige}\end{blocLivre}
\begin{blocLivre}{cor-7-4}\begin{corrige}{4}{Même dénominateur}\demarche{À parts égales, comparer les numérateurs.}
Les huitièmes ont la même taille : $\res{\frac58>\frac38}$. Pour les cinquièmes : $\res{\frac15<\frac45<\frac75}$.
\end{corrige}\end{blocLivre}
\begin{blocLivre}{cor-7-5}\begin{corrige}{5}{Une part d’une quantité}\demarche{Partager la quantité selon le dénominateur.}
Partager 12 en 3 : $12:3=\res4$ billes. Partager 100 en 4 : $100:4=\res{25}$ mètres.
\end{corrige}\end{blocLivre}
\begin{blocLivre}{cor-7-6}\begin{corrige}{6}{Réunir des parts}\demarche{Ajouter les numérateurs, garder le dénominateur.}
Les parts s’ajoutent : $\res{\frac58}$. Ensuite $\frac25+\frac35=\frac55=\res1$ : l’unité est reconstituée.
\end{corrige}\end{blocLivre}
\begin{blocLivre}{cor-7-7}\begin{corrige}{7}{Retirer des parts}\demarche{Écrire 1 avec le dénominateur utile.}
$\res{\frac7{10}-\frac4{10}=\frac3{10}}$. Pour la seconde, écrire $1=\frac99$ : $\frac99-\frac29=\res{\frac79}$.
\end{corrige}\end{blocLivre}
\begin{blocLivre}{cor-7-8}\begin{corrige}{8}{Compléter jusqu’à 1}\demarche{Compter les parts manquantes.}
Il faut 9 neuvièmes en tout : $\square=\res4$. Pour les sixièmes : $\square=\res5$, car $\frac56+\frac16=\frac66=1$.
\end{corrige}\end{blocLivre}
\begin{blocLivre}{cor-7-9}\begin{corrige}{9}{Des cinquièmes bien placés}\demarche{Convertir l’écriture mixte en pas sur la droite.}
E est 4 pas après 0 : $\res{E\left(\frac45\right)}$. F se place 2 pas après 1, c’est $\frac75$ ; sa position figure en vert.
\label{fig-cor-7-9}\figureGeo[corrige]{fd-9}
\end{corrige}\end{blocLivre}
\begin{blocLivre}{cor-7-10}\begin{corrige}{10}{Quelle part est la plus grande ?}\demarche{À numérateur égal, le plus petit dénominateur gagne.}
Partager en 3 donne des parts plus grandes que partager en 4 : $\res{\frac13>\frac14}$. En prenant deux parts de chaque : $\res{\frac23>\frac24}$. Les bandes superposées le montrent sans calcul.
\label{fig-cor-7-10}\figureGeo[corrige]{fd-10}
\end{corrige}\end{blocLivre}
\begin{blocLivre}{cor-7-11}\begin{corrige}{11}{Mesurer avec une bande}\demarche{Reporter l’unité entière avant les parts.}
La bande se reporte 2 fois en entier, puis 3 plis d’un quart complètent : la longueur vaut $\res{2+\frac34}$ d’unité, soit $\frac{11}4$.
\label{fig-cor-7-11}\figureGeo[corrige]{fd-11}
\end{corrige}\end{blocLivre}
\begin{blocLivre}{cor-7-12}\begin{corrige}{12}{La randonnée}\demarche{Calculer la part puis le reste du tout.}
Le quart de 20 km : $20:4=\res5$ km le matin. Il reste $20-5=\res{15}$ km. Contrôle : $5+15=20$ km, le sentier entier.
\end{corrige}\end{blocLivre}
\begin{blocLivre}{cor-7-13}\begin{corrige}{13}{Sept parts de deux tartes}\demarche{Écrire les deux tartes en sixièmes avant de soustraire.}
Deux tartes font $\frac{12}6$. Il reste $\frac{12}6-\frac76=\res{\frac56}$ de tarte : un peu moins qu’une tarte entière.
\end{corrige}\end{blocLivre}
\begin{blocLivre}{cor-7-14}\begin{corrige}{14}{La même longueur ?}\demarche{Superposer les parts pour comparer.}
$\res{\text{Oui}}$ : deux quarts superposés couvrent exactement une moitié, donc $\frac12=\frac24$. Deux écritures différentes peuvent désigner le même nombre.
\label{fig-cor-7-14}\figureGeo[corrige]{fd-14}
\end{corrige}\end{blocLivre}
\begin{blocLivre}{cor-7-15}\begin{corrige}{15}{Le ruban partagé}\demarche{Traiter chaque part unitaire comme un partage.}
Ana : $18:2=\res9$ m. Bilal : $18:3=\res6$ m. Ensemble ils prennent 15 m ; Chloé reçoit $18-15=\res3$ m. Contrôle : $9+6+3=18$ m.
\end{corrige}\end{blocLivre}
